\section{Conclusion}

We formulated context selection through the set-conditioned marginal reduction
in end-task loss. This decision object exposes structural headroom created by
state conditioning, and it separates that opportunity from the accuracy with
which a router can estimate it. R-MUR combines cached screening with compact
predictor-response summaries and a gap-weighted ranking objective. Across six
standard traffic benchmarks, R-MUR improves learned context selection over
standalone and cached baselines, and a multi-target study shows that structural
headroom persists across sensors.

Two additions sharpen the scope of that claim. First, a general smooth-loss
analysis shows that marginal utility decomposes into the alignment between the
predictor's response and its residual plus a curvature penalty, with
cross-entropy and other Bregman losses covered by exact or two-sided
statements; predictor response is therefore the theoretically matched
observation, and the numerical study confirms the bound empirically. Second, a
subset-capable nonlinear backbone preserves structural headroom but reduces the
share of it that a router can realize from observable features: the learned
response-free routers become indistinguishable from random selection and
R-MUR's advantage over them is no longer statistically resolved, while a second
task family whose candidates are not redundant is already solved by
nearest-neighbour relevance. Predictor response and decision margin remain the
key quantities for turning structural opportunity into learned routing gains,
and the strength of the fixed predictor determines how much of that opportunity
is identifiable before the decision is made.
